\documentclass[10pt,a4paper]{article}
\usepackage[T1]{fontenc}
\usepackage{lmodern}
\usepackage[margin=23mm]{geometry}
\usepackage{amsmath,amssymb,booktabs,tabularx,microtype}
\usepackage[hidelinks,hypertexnames=false]{hyperref}
\pdfinfoomitdate=1
\pdftrailerid{}
\pdfsuppressptexinfo=15
\setlength{\parindent}{0pt}
\setlength{\parskip}{5pt}
\newcommand{\kb}{k_{\mathrm B}}
\begin{document}
\begin{center}
{\Large A native atmosphere with composition feedback}\\[4pt]
Definitions and validation obligations for the selected solar proxy
\end{center}

\section{What is solved}
\path{escape.solar.SolarCompositionModel} solves a common radiation field,
atomic populations and a reduced chemical network at supplied temperature
and density. \path{SolarAtmosphereModel.solve_atmosphere} additionally closes
temperature, molecular weight, mass conservation and transonic momentum.
The solved species supply the observer's continuum and H$\alpha$ line.
Cloudy and MCEG are external verification tools, never runtime dependencies.

The target composition is the user's selected solar proxy, using the stated
Cloudy C23 default abundances. It is not a recovered author input file.
H and He are coupled to C, N, O, Mg, Al, Si, Ca, Mn and Fe at these abundances;
omitted elements and processes need an explicit contribution bound.
The retained metal ladders contain charges $q=0,1,2$. Resolved excitation
atoms are C I/II, N I, O I, Mg II, Al II, Si II, Ca II, Mn II and Fe II/III.
The current Fe II candidate retains the available 716 levels; 80 and 371
remain explicit atomic-convergence controls. The 80-level reference
substantially underestimated UV pumping in the tested solar columns.
Helium uses a three-charge-state ground-atom reduction. This definition
does not itself establish adequacy for a solar atmosphere.

\section{Conservation and pressure}
Let $n_{\rm H}$ count hydrogen nuclei, $n_i$ the bound atomic H populations,
and $n_-,n_{2+},n_2^{\rm mol},n_{2*}^{\rm mol}$ the densities of H$^-$,
H$_2^+$ and the two H$_2$ superstates. Then
\begin{align}
 n_{\rm H}&=\sum_i n_i+n_p+n_-+2(n_{2+}+n_2^{\rm mol}+n_{2*}^{\rm mol}),
 \label{eq:nuclei}\\
 n_e&=n_p-n_-+n_{2+}+\sum_q q n_{{\rm He},q}
            +\sum_{Z,q}q n_{Z,q},\qquad
 \sum_q n_{Z,q}=A_Z n_{\rm H}. \label{eq:charge}
\end{align}
Softmax coordinates conserve each element; a molecular hydrogen fraction
is divided by its number of H nuclei to recover its particle density.
Electrons are calculated from charge neutrality, including H$^-$.
With $n_{\rm part}$ counting all material particles including free electrons,
\begin{equation}
 \rho=n_{\rm H}m_{\rm per\,H},\qquad
 \mu=\frac{\rho}{m_p n_{\rm part}},\qquad
 P=n_{\rm part}\kb T,\qquad a^2=\frac{\kb T}{\mu m_p}.
 \label{eq:pressure}
\end{equation}
The mass per H nucleus includes the retained elemental abundances, rather
than imposing the pure-H relation $\mu=(1+n_p/n_{\rm H})^{-1}$.

\section{Radiation and populations form a single problem}
The continuum coefficients are sums over the actual solved species:
\begin{equation}
 \chi_\nu=\sum_s\chi_{\nu,s},\quad j_\nu=\sum_s j_{\nu,s},\quad
 \mu_\Omega\frac{dI_\nu}{dr}=j_\nu-\chi_\nu I_\nu,\quad
 J_\nu=\tfrac12\int_{-1}^{1}I_\nu\,d\mu_\Omega.
 \label{eq:rt}
\end{equation}
The implementation uses exact homogeneous-cell integration. A normal stellar
beam and diffuse emission are distinguished. Thomson/Rayleigh scattering
use the declared attenuation approximation; they are not counted as thermal
absorption. H bound--free, molecular continua, metal/He continua, shared
free--free and spontaneous two-photon emission enter the common solve.
The Gaunt factors follow the C23 tables~\cite{gaunt1,gaunt2}.

Lines use an explicit escape/destruction closure instead of adding escape
probabilities to a full trapped-line radiation field. A representative pair is
\begin{equation}
 R_{lu}=B_{lu}\overline J_{\rm pump},\qquad
 R_{ul}=A_{ul}(\beta_{\rm esc}+\beta_{\rm dest})
                +B_{ul}\overline J_{\rm pump}.
 \label{eq:line}
\end{equation}
The specialized H Ly$\alpha$ and metal prescriptions are documented in
\path{escape_probability.tex} and the source. Direct stellar pumping is
evaluated at each actual line center: interpolating a coarse continuum grid
can miss a narrow stellar absorption trough. Diffuse spectra may be
interpolated separately. The selected static-column reference does not shift
the incident stellar pumping spectrum into each flowing cell's rest frame;
velocities enter the wind, energy terms and observer. This matches the stated
public-sunbather-style comparison, not a moving-medium line-transfer claim.
The solar reference explicitly doubles the total
line column, following its \texttt{double optical depths} input; the finite
open-slab control uses an undoubled column. Neither setting is a fit parameter.

H Ly$\beta$ and the six overlapping O I components instead share a resolved
CRD transfer calculation. Their isolated line rates and old destroyed-photon
sources are replaced together. The change in radiation seen by continuum
absorbers is
\begin{equation}
 \delta\Phi_\nu=\frac{4\pi\,\Delta\nu}{h\nu}
 (J_\nu^{\rm H+O+cont}-J_\nu^{\rm cont}).\label{eq:overlap}
\end{equation}
This signed correction includes both extra line photons and screening of
continuum photons. Every retained species receives the same correction.
The group uses finite slab boundaries and actual line centers; it is an
explicit alternative to C23's inconsistent auxiliary O I pumping formula
(\path{species2.cpp:setXtraRatesO1}, ticket 78).
Its LTE, photon-budget and derivative checks do not alone establish
the atmospheric impact of choosing CRD rather than partial redistribution.

Locally destroyed line photons are allocated to all continuum absorbers
using the same opacity denominator. If $D_\ell$ is their volume production,
\begin{equation}
 \Phi^{\rm OTS}_\ell=\frac{D_\ell}{\chi_{\nu_\ell}},\qquad
 \Gamma_{s,\ell}=\sigma_s(\nu_\ell)\Phi^{\rm OTS}_\ell.
 \label{eq:ots}
\end{equation}
The corresponding excess photoelectron energy is counted at the receiving
species, once. Escape itself is not an additional heating term.
Atomic and chemical populations obey
\begin{equation}
 0=\sum_j(n_jR_{ji}-n_iR_{ij})+S_i^{\rm ion/rec/chem},\qquad
 0=\sum_r\nu_{sr}\,k_r(T,J_\nu)\prod_a n_a^{\alpha_{ar}}.
 \label{eq:rates}
\end{equation}
Thus a change in ionization changes opacity, electron collisions, molecular
abundances and pressure, which change the radiation and flow in return.

\section{Energy and wind}
The kinetic ledger counts collisional level-energy exchange, photoelectron
excess energy, recombination, free--free and specified molecular processes:
\begin{align}
 Q_{\rm ph,s}&=n_s\int \frac{4\pi J_\nu}{h\nu}\sigma_s(\nu)
                    (h\nu-\chi_s)\,d\nu,\label{eq:photoheat}\\
 0&=\Gamma-\Lambda-\rho v\frac{d}{dr}
       \left(\frac{3\kb T}{2\mu m_p}\right)
       +\frac{\kb T v}{\mu m_p}\frac{d\rho}{dr}.\label{eq:energy}
\end{align}
The last two terms are the public sunbather local-SE thermal-flow closure
\cite{sunbather}. Population advection and advected excitation/ionization
energy are absent. The supplied innermost temperature replaces one energy
equation unless the explicit dense lower patch below is selected.
An explicitly selected radiative-equilibrium comparison omits the
two flow-energy terms while retaining momentum feedback.

\paragraph{Dense lower boundary.}
Zhang et al. exclude levels above $10^{-3}$ bar from the Cloudy thermal
calculation and extrapolate toward a 4500 K photosphere~\cite{zhang}.
The optional \texttt{dense\_lower\_cells} selects a fixed lower subset;
its uppermost join must leave the non-LTE cells below $10^{-3}$ bar and
$10^{15}\,\mathrm{cm}^{-3}$. The conditional native prescription is
\begin{equation}
 T(r)=T_{\rm base}+\frac{r-R_p}{r_j-R_p}(T_j-T_{\rm base}),\qquad
 \frac{n_q}{n_{q+1}}=\frac{n_e S}{2T^{3/2}}
       \frac{U_q(T)}{U_{q+1}(T)}\exp\!\left(\frac{\chi_q}{\kb T}\right),
 \label{eq:lower}
\end{equation}
where $r_j$ is the first non-LTE cell center and $S$ the Saha constant.
Atomic excitation and ionization use the actual conserved electron density
and retained-level partition functions $U_q$; unrepresented charge-state
excitation uses the packaged bound-level partition data. Ground-state helium
retains its declared reduction. H$^-$ obeys its Saha relation; the other
three molecular species retain their reaction equations. Dense metal
continuum emission is replaced by $j_\nu=\chi_\nu B_\nu(T)$, including removal
of the independent recombination pseudo-line sources there. Combining full
ionic partitions with the ground-only Milne prescription would otherwise
violate Kirchhoff's law. The non-LTE upper prescription is unchanged.
These cells supply opacity, radiation and pressure but do not claim
non-LTE atomic SE or thermal equilibrium. The wind includes an explicit
$R_p,T_{\rm base}$ knot with the innermost solved $\mu$.
The author join law and dense populations are unavailable: this boundary
must be varied and its spectral effect bounded. A converged numerical root
outside the stated upper-atmosphere domain is not an applicable spectrum.

In non-LTE cells, H$^-$ obeys the reaction network. H$_2^+$ formation and
H$^-$ attachment retain C23's empirical cooling, which need not match the
formation-rate fit. Generic reaction enthalpies must not be added again.
The two-superstate H$_2$ and ground-He reductions require contribution tests.

The non-LTE heavy-ion ground Milne and two RR+DR cascade photons retain
C23's independent prescriptions: the cascade does not subtract ground
capture. This approximate source budget requires a sensitivity check;
clipping a negative ground-subtracted remainder would hide the issue.

Mass conservation and momentum close the atmosphere:
\begin{equation}
 \rho=\frac{\dot M}{4\pi r^2v},\qquad
 \left(v-\frac{a^2}{v}\right)\frac{dv}{dr}
 =\frac{2a^2}{r}-\frac{da^2}{dr}-\frac{GM_p}{r^2}
                      +\frac{3GM_\star}{a_{\rm orb}^3}r.
 \label{eq:wind}
\end{equation}
The regular sonic branch comes from the existing nonisothermal Parker
solver. The sampled $T/\mu$ is continued constantly to $32R_p$; that outer
extension is an approximation to test. It is not a computed outer thermal
profile. The linear stellar tide is an effective gravity term, not a 3D
Roche surface.

\section{Derivatives and scientific checks}
For the complete residual $F(y,\theta)=0$, including radiation, chemistry,
temperature and wind,
\begin{equation}
 \frac{dy}{d\theta}=-F_y^{-1}F_\theta.\label{eq:implicit}
\end{equation}
Newton--GMRES applies exact JAX Jacobian products without a dense global
Jacobian. The implicit tangent and transpose solves check their residuals.
Unconverged roots or derivatives are not valid spectra for inference.
At exactly zero FUV (including floating-point transfer underflow), the
empirical small-H$_2$ $J^{1/5}$ term uses a zero tangent; its positive-field
formula is unchanged. This is an explicit dark-boundary convention, not a
claim of a finite right derivative with respect to $J$ at zero.
The observer consumes these same populations and continua on the same
radial edges; rotation, zonal velocity and radial velocity remain explicit.

MCEG separates component tests, complete-root/derivative checks, reference
agreement and observational applicability. C23's normal stopping rules do
not guarantee converged populations or total optical depths; the reference
requires refinement tests too. Other overlap channels, excited-He effects,
omitted species, finite H shells, lower boundaries, numerical grids and
observation processing remain requirements until evidence resolves them.

\begin{thebibliography}{9}
\bibitem{zhang} Zhang et al. (2026),
\href{https://arxiv.org/abs/2606.23899v1}{arXiv:2606.23899v1}.
\bibitem{sunbather} Linssen et al. (2024),
\href{https://arxiv.org/abs/2404.12775}{arXiv:2404.12775}; public sunbather
energy and hydrodynamic closure. Author modifications are not recovered.
\bibitem{cloudy} Chatzikos et al. (2023), RMxAA 59, 327; Cloudy C23.01
source and its distributed atomic/chemical data. Altered JAX adaptations
retain the original license and extraction provenance.
\bibitem{gaunt1} van Hoof et al. (2014), MNRAS 444, 420.
\bibitem{gaunt2} van Hoof et al. (2015), MNRAS 449, 2112.
\end{thebibliography}
\end{document}
